
\subsubsection{Direct Preference Optimization}

Direct Preference Optimization (DPO) provides a closed-form solution for learning from human preferences without training a separate reward model \citep{rafailov2023direct}. Given paired preferences (y\_w, y\_l) for prompt x, DPO optimizes:

$$\mathcal{L}_{\text{DPO}} = -\log \sigma\left(\beta \log \frac{\pi_\theta(y_w|x)}{\pi_{\text{ref}}(y_w|x)} - \beta \log \frac{\pi_\theta(y_l|x)}{\pi_{\text{ref}}(y_l|x)}\right)$$

This formulation has become a dominant approach for preference-based fine-tuning. However, DPO optimizes a single objective: matching human preferences. The present work asks whether auxiliary objectives can capture dimensions of response quality orthogonal to preference labels.

\subsubsection{Multi-Objective Extensions to DPO}

Several works extend DPO with additional objectives. MODPO introduces margin-based auxiliary losses that provide additional training signal beyond binary preferences \citep{zhou2024modpo}. PAMA frames alignment as multi-objective optimization across multiple preference dimensions \citep{he2025pama}. These methods share a common pattern: they add objectives derived from preference data. BiDPO differs by adding an objective orthogonal to preference labels.

\subsubsection{Human Agency in AI Systems}

Mitelut et al. (2023) argue that intent-aligned AI may inadvertently deplete human agency by optimizing for immediate task completion rather than capability transfer. This theoretical framework motivates the present empirical investigation into training objectives that preserve agency-related patterns.

